\section{Methodology}
\subsection{\dslr{} Formulation}
We describe \dslr{}'s architecture here - generator, discriminator, contrastive learning for the discriminator and the training setup for MOVES, below.
\subsubsection{Generator}
We use range image-based LiDAR preprocessing, which allows us to use image-based processing on LiDAR effectively. We use a conditional GAN architecture for the  static translation. We use the Encoder ($E_\phi$) and Decoder ($D_\theta$) from (\citet{caccia2018deep}) for our generator (\textit{G}). Given $d_i\in \textbf{Dy}$ as input, the goal of $G$ is to generate $\overline{s_i}$  (\textit{s.t.} $\overline{s_i}$ resembles $s_i$). 
\subsubsection{Discriminator}
\label{subsub-disc}
We provide an overview of the challenges while designing the discriminator. \textbf{(1)} Discriminating a static scan ($s_i$) from the dynamic scan ($d_i$) is challenging because both are generated from similar LiDAR distribution and share similar patterns - scenes and structure. Dynamic LiDAR scans are equivalent to static LiDAR with noise in the form of dynamic objects. Latent feature vectors of both consists of mostly homogeneous and almost indistinguishable patterns. Features and patterns that help in discrimination (i.e. dynamic objects and occlusions) are limited because they occupy a small percentage of points in the scan compared to the common static regions. This leads to an imbalance w.r.t. the discriminatory features, making them difficult to learn. \textbf{(2)} A vanilla discriminator for static scans \textit{v/s} dynamic scans, suffers from a \textit{simplicity bias} (\citet{shah2020pitfalls}) - the discriminator needs to only pay attention to the \textit{simplest features} for discrimination i.e. (moving objects), to distinguish them from static scans. It saturates prematurely by learning these simplest features. Movable objects (which are difficult to discover due to uncertain motion) do not contribute significantly to the learning because the model never needs to use these complex movable object features. This leads to  sub-optimal translations as shown in Ablation studies (using a GAN with a vanilla discriminator) (Section \ref{sec:ablation} and Table \ref{tab:ablation-disc}).

To address the above challenges, we design a novel \textit{couple discriminator} ($D$) for our proposed GAN, which operates on the latent space of the LiDAR scans. $D$ takes a pair of LiDAR latent representations as input.

Let ($r_{s_i}$, $r_{s_j}$, $r_{d_k}$) be the latent representation for the triplet (${s_i}$, ${s_j}$, ${d_k}$) where $s_i$, $s_j \in \textbf{St}$ and $d_k \in \textbf{Dy}$. We construct 2 pairs from the triplet: ($r_{s_j}$, $r_{s_i}$) which is labelled as real, and ($r_{d_k}$, $r_{s_i}$), which is labelled as fake by the \textit{D}.

In order to classify the above pairs as real and fake respectively, \textit{D} is forced to  amplify discriminatory features in dynamic scan, $d_k$. Instead of competing static representations against dynamic (as in a vanilla discriminator), the discrimination occurs between similarities in static scans in the  pair ($r_{s_j}$, $r_{s_i}$) v/s differences between static and dynamic scan in the pair ($r_{d_k}$, $r_{s_i}$). This helps $D$ to focus only on the discriminatory features between static and dynamic scans without using segmentation labels. $D$ can explicitly focus on dynamic objects and occlusions in $d_k$ while ignoring the common static regions. It is also able to differentiate between occlusions due to static objects and occlusions due to dynamic objects, thereby enabling \textit{G} to only remove dynamic occlusion and objects while retaining static occlusions and static objects. This behaviour is further assisted by the absence of moving and movable objects in static scans.

Further, to stabilize the training of $D$, we replace the negative log-likelihood objective of the discriminator with the least square objective (\citet{mao2017least}). The idea behind using Least Square GANs is two fold: first is to ensure that $D$ receives smooth and non-saturating gradients, thereby allowing $G$, which is also trained via gradients through $D$, to be trained effectively. Secondly, given gradients could be weak due to the structural and pattern similarity between static and dynamic scans, we would even want to penalize LiDAR pairs that are far from the discriminator decision boundary,  though on the correct side. This allows $G$ to generate static scans closer to the decision boundary to avoid learning saturation (\citet{mao2017least}). It helps to generate more gradients while updating $G$, which allows reconstructed samples to be closer to the real manifold of static scans.

\textbf{Contrastive Learning on  Discriminator:} Dynamic and static scans are generated from similar distributions and share similar patterns (except for movable and moving objects). Learning discriminatory features between them is non-trivial. To assist $D$ in paying attention only to the discriminatory patterns between them, we use a contrastive loss to \textit{contrast dynamic scans from static} ones while bringing static scans closer. Here we use a careful choice of static and dynamic scans using \textit{hard negative mining}. Recall the triplet above $(s_i\in St, s_j\in St, d_k\in Dy)$, used while training the discriminator. In this triplet, we now restrict \textit{$d_k$ to be the corresponding dynamic scan for $s_{i}$} ($i.e.$ $d_{k}$ is now $d_{i}$).
\begin{figure}[t]
  \centering
  \includegraphics[width=0.6\linewidth, height=0.5\linewidth]{ECCV-Final.pdf}  
  \caption{Architecture Diagram for \dslr{}}
\label{fig:carlagan}
\end{figure}
We now explain the reason behind this choice. We use the notation (${s_i}$, ${s_j}$, ${d_k}$) instead of ($r_{s_i}$, $r_{s_j}$, $r_{d_k}$) for the sake of simplicity.
Consider $s_i^+$ as the anchor, $s_j^+$ as positive, and $d_i^-$ \textit{(corresponding dynamic/hard negative for $s_i^+$)} as the negative scan. In the ideal case, the anchor $s_i^+$, and the positive scan $s_j^+$ should be closer, while the anchor and the negative scan should be far away in the latent space. However due to our careful selection of the   negative sample, the anchor $s_i^+$, is more closer to the negative sample $d_i^-$, than the positive sample $s_j^+$ in the latent space. This is because the anchor $s_i^+$, and the negative sample $d_i^-$ are a corresponding pair and have very high similarity despite belonging to different classes.


The samples $s_i^+$ and $d_i^-$ are a corresponding static-dynamic pair, are perceived in the same environment and share the same static structures and patterns. Thus, they are very close in the latent space. The only difference between them - movable and moving objects, accounts for the discriminatory features. Contrasting them \textit{w.r.t.} each other forces $D$ to implicitly focus on the only available differences between them - moving and movable objects in $d_i^-$, and the corresponding static background features in static scan ($s_i^+$). This ensures that the signals that flow to $G$, explicitly focus on the above discriminatory patterns and structures, thereby helping to identify moving and movable objects and corresponding occlusions.
It also helps in differentiating \textit{occlusion due to dynamic objects from occlusions due to static objects}, thereby overcoming limitations of a vanilla discriminator as discussed in the previous sections. These improvements strengthen the discriminatory signals to $D$ while also helping $G$ focus on dynamic objects and occlusions explicitly, for reconstructing background static scenes. Given \textit{m} is a distance metric, and $\alpha$ is the neighbourhood parameter, contrastive triplet loss, $L_{tri}$ is defined as:
\begin{equation}
        =max(0, m( s_i^+, s_j^+)  - m(d_i^-, s_i^+) + \alpha)
\end{equation}


\subsection{Training Setup for Conditional Translation}
\label{setup-carla}


 \textit{Discriminator}: We use a MLP-based sequential neural network consisting of 6 layers with sigmoid non-linearity. As detailed above, for training, we consider a triplet of the form $(s_i\, s_j, d_i)$ (Figure \ref{fig:carlagan}) for training $D$. We generate latent vectors of the triplet using the encoder ($E_{\phi}$ of $G$), and create pairs ($r_{s_j}$, $r_{s_i}$) and ($r_{d_i}$, $r_{s_i}$) using the triplet. The former is labeled as real while the latter as fake by our discriminator. $D$ is trained using the discriminatory signals from the real $v/s$ fake classification and  the Contrastive triplet loss.

\textit{Generator:} The generator $G$ is fed with a dynamic scan ($d_i$) and is trained with the reconstruction loss using the corresponding static scan $s_i$ (Figure \ref{fig:carlagan}). Note that the reconstructed scan does not flow into the discriminator $D$. Its latent vector $r_{d_i}$ flows in the discriminator because we operate on the latent space. This latent vector ($r_{d_i}$) is a constituent of the pair ($r_{d_i}$, $r_{s_i}$), which is labeled as fake by $D$ (Figure \ref{fig:carlagan}).
The generator is trained using the gradients received from $D$, as well the gradient signals from the decoder($D_{\theta}$ of $G$).


To enhance the gradients to $G$ and  $D$, we employ a relativistic least square GAN for \dslr{}.
    The generator $G(\phi,\theta{})$ takes $d_i\in Dy$ as input and generates $\overline{s_i} \sim s_i$ (where ${s_i} \in St$  is the corresponding static of $d_i$). On the other hand, given a triplet of LiDAR scan latent representations $(r_{s_i}, r_{s_j}, r_{d_i})$, the couple discriminator ($D$) classifies the couple $(r_{s_j}, r_{s_i})$ as real and $(r_{d_i}, r_{s_i})$ as fake. 

     We list the loss formulation for the Discriminator ($Loss_{D}$) as 
    \begin{equation}
    \begin{split}
           &= D((r_{s_j},r_{s_i}), real) + D((r_{d_i},r_{s_i}), fake)\\ &+ {L_{tri}}(r_{s_i}, r_{s_j}, r_{d_i})
    \end{split}
    \label{eq:disc}
    \end{equation}
    and for the generator ($Loss_{G}$) as 
    \begin{equation} 
    \begin{split}   
    &= D((r_{d_i},r_{s_i}), real) +Recon(\overline{d_i},s_i)
    \end{split}
    \label{eq:gen}
    \end{equation}
    where \textit{Recon}() denotes the Reconstruction Error while $L_{tri}$ denote the Contrastive triplet loss and $\overline{d_i}$ is the reconstructed output of the $G$, given $d_i$ as input.

    

    Assuming $P(St)$ denotes the distribution for the static LiDAR scans, while $P(Dy)$ denotes the distribution for the dynamic LiDAR scans. \\

    
    Equation \ref{eq:disc} can be expanded as
    \begin{equation}
    \begin{split}
     & =
    {\mathbb{E}}_{s_i,s_j\sim P(St)}[(DI(r_{s_j},r_{s_i}) -1)^{2}]\\
    & + {\mathbb{E}}_{d_i\sim P(Dy),s_i\sim P(St)}[(DI(r_{d_i},r_{s_i}) - 0)^{2}] \\ 
    & + {L_{tri}}(r_{s_i}, r_{s_j}, r_{d_i})
    \end{split}
    \end{equation}
    while equation \ref{eq:gen} can be expanded as
    \begin{equation}
    \begin{split}
         =
         {\mathbb{E}}_{d_i\sim P(Dy),s_i\sim P(St)}[(DI(r_{d_i},r_{s_i}) -1)^{2}] +
    Recon(\overline{d_i}, s_i)
    % }
    \end{split}
    \label{egenloss}
    \end{equation}







Now we explain the training process of \dslr{}.
Given a batch of triplets $(s_i,s_j,d_i)$, we use Encoder ($E_{\phi}$) of $G$, to extract the latent representations, i.e. $(r_{s_i},r_{s_j},r_{d_i})$. 
We create the real and fake pairs, ($r_{s_j}$, $r_{s_i}$) and ($r_{d_i}$, $r_{s_i}$) respectively. $D$ it is trained to classify these pairs correctly during training. $G$ is trained
% given a dynamic scan ($d_i$) as input, 
to produces a latent representation ($r_{d_i}$), such that the pair ($r_{d_i}$, $r_{s_i}$) is labeled as real instead of fake by $D$. In order to achieve this, \textit{G} must ensure that for the given dynamic scan $d_i$, it must produces a latent representations that is close to the corresponding static LiDAR scan $r_{s_i}$, on the static manifold. This allows us to achieve our translation.
 
\begin{figure}[t]
  \centering
  \includegraphics[width=0.6\linewidth, height=0.6\linewidth]{ECCV-Kitti-Final.pdf}  
  \caption{Architecture Diagram for \dslruda{}}
\label{fig:ARCHDAG}
\end{figure}
\subsection{\dslruda{} }
A challenge for our approach is the unavailability of corresponding ground truth static for dynamic LiDAR scans for some datasets. To overcome this, we resort to Domain Adaptation and  follow (\citet{Borgwardt2006IntegratingSB}) to integrate a domain discrepancy loss ($L_{domain\_distance}$) between the paired and the unpaired domain in the latent space. There exists several other methods and loss functions that can be used here. We use (\citet{Borgwardt2006IntegratingSB}) for adaptation because it is straight forward, requires minimal change to integrate in \dslr{} and works well. It projects the data points from both domains in the reproducible Kernel Hilbert space followed by computing the average difference between them. For our purpose the source domain is CARLA-64 dataset (\citet{dosovitskiy2017carla}) and the target domain is KITTI dataset (\citet{geiger2012we}). We consider the same architecture of \dslr{} for \dslruda{} with a few additions (Figure \ref{fig:ARCHDAG}) -
\\ \textbf{(a)} We introduce a new generator for KITTI - $G_{K}(E_{\gamma}, D_{\lambda})$. KITTI scans are labelled as $k_l$.
\\ \textbf{(b)} Latent representations of KITTI scan ($E_{\gamma}(k_l)$ or $r_{k_l}$) is paired with $r_{s_j}$ :  ($r_{k_l}, r_{s_j}$). It is labelled as fake by $D$. \\
 \textbf{(c)} Domain distance loss is calculated between $r_{k_l}$ (latent vector of dynamic scan in target domain) and $r_{s_i}$ (latent vector of static scan in source domain).\\ 
 \textbf{(d)} KITTI dynamic scan is reconstructed using the KITTI generator $G_{K}(E_{\gamma}, D_{\lambda})$. It helps to retain the static structures of the target scan. 

We use a triplet of latent representations obtained from LiDAR scans in the source CARLA-64 domain $(r_{s_i}, r_{s_j}, r_{d_i})$ as discussed in Section \ref{setup-carla}. We use another latent representation derived from the target KITTI domain, $r_{k_l}$  using the KITTI generator $G_{K}(E_\gamma,D_\lambda)$.


    Discriminator Loss for \dslruda{}, $Loss_{D-kitti}$ is as follows:
    \begin{equation}
    \begin{split}
           &= Loss_{D} +  D((r_{k_l}, r_{s_j}), fake)  
    \end{split}
    \end{equation}  
    \begin{equation}
        \begin{split}        
           =Loss_{D} +  {\mathbb{E}}_{k_l\sim P(kitti),s_j\sim P(St)}[(D(r_{k_l},r_{s_j}) - 0)^{2}] 
    \end{split}
    % \label{eq:disc}
    \end{equation}
    and the generator loss $Loss_{G-kitti}$ is as follows:
    \begin{equation}
    \begin{split}
           &= Loss_{G} +  Recon(\overline{k_l}, k_l)
           +L_{Domain-distance}(r_{k_l},r_{s_j})
    \end{split}
    % \label{eq:disc}
    \end{equation}